\textbf{Keywords:} Systems, Feedback, PLL, Integer Divider, SD, SD PLL,
Modulation, linear phase model

\section{Why clocks?}\label{why-clocks}\index{why clocks?@Why clocks?}

Virtually all integrated circuits have some form of clock system.

For digital we need clocks to tell us when the data is correct. For
Radio's we need clocks to generate the carrier wave. For analog we need
clocks for switched regulators, ADCs, accurate delay's or indeed, long
delays.

The principle of a clock is simple. Make a 1-bit digital signal that
toggles with a period \(T\) and a frequency \(f = 1/T\).

The implementation is not necessarily simple.

The key parameters of a clock are the frequency of the fundamental,
noise of the frequency spectrum, and stability over process and
enviromental conditions.

When I start a design process, I want to know why, how, what (and
sometimes who). If I understand the problem from first principles it's
more likely that the design will be suitable.

But proving that something is suitable, or indeed optimal, is not easy
in the world of analog design. Analog design is similar to physics. An
hypothesis is almost impossible to prove ``correct'', but easier to
prove wrong.

\subsection{A customer story}\label{a-customer-story}\index{a customer story@A customer story}

Take an example.

\subsubsection{Imagine a world}\label{imagine-a-world}

\begin{quote}
``I have a customer that needs an accurate clock to count seconds''. --
Some manager that talked to a customer, but don't understand details.
\end{quote}

As a designer, I might latch on to the word ``accurate clock'', and
translate into ``most accurate clock in the world'', then I'd google
atomic clocks, like
\href{https://en.wikipedia.org/wiki/Rubidium_standard}{Rubidium
standard} that I know is based on the hyperfine transition of electrons
between two energy levels in rubidium-87.

I know from quantum mechanics that the hyperfine transition between two
energy levels will produce an precise frequency, as the frequency of the
photons transmitted is defined by \(E = \hbar \omega = h f\).

I also know that quantum electro dynamics is the most precise theory in
physics, so we know what's going on.

As long as the Rubidium crystal is clean (few energy states in the
vicinity of the hyperfine transition), the distance between atoms stay
constant, the temperature does not drift too much, then the frequency
will be precise. So I buy a
\href{https://www2.mouser.com/ProductDetail/IQD/LFRBXO059244Bulk?qs=iw0hurA\%2FaD0K8weKx\%2Fu2ow\%3D\%3D}{rubidium
oscillator} at a cost of \$ 3k.

I design an ASIC to count the clock ticks, package it plastic, make a
box, and give my manager.

Who will most likely say something like

\begin{quote}
``Are you insane? The customer wants to put the clock on a wristband,
and make millions. We can't have a cost of \$ 3k per device. You must
make it smaller and it must cost 10 cents to make''
\end{quote}

Where I would respond.

\begin{quote}
``What you're asking is physically impossible. We can't make the device
that cheap, or that small. Nobody can do that.''
\end{quote}

And both my manager and I would be correct.

\subsubsection{Imagine a better world}\label{imagine-a-better-world}

Most people in this world have no idea how things work. Very few people
are able to understand the full stack. Everyone of us must simplify what
we know to some extent. As such, as a circuit designer, it's your
responsibility to fully understand what is asked of you.

When someone says

\begin{quote}
'' I have a customer that needs an accurate clock to count seconds''
\end{quote}

Your response should be ``Why does the customer need an accurate clock?
How accurate? What is the customer going to use the clock for?''. Unless
you understand the details of the problem, then your design will be
sub-optimal. It might be a great clock source, but it will be useless
for solving the problem.

\subsection{Frequency}\label{frequency}\index{frequency@Frequency}

The frequency of the clock is the frequency of the fundamental. If it's
a digital clock (1-bit) with 50 \% duty-cycle, then we know that a
digital pulse train is an infinite sum of odd harmonics, where the
fundamental is given by the period of the train.

\subsection{Noise}\label{noise}\index{noise@Noise}

Clock noise have many names. Cycle-to-cycle jitter is how the period
changes with time. Jitter may also mean how the period right now will
change in the future, so a time-domain change in the amount of
cycle-to-cycle jitter. Phase noise is how the period changes as a
function of time scales. For example, a clock might have fast period
changes over short time spans, but if we average over a year, the period
is stable.

What type of noise you care about depends on the problem. Digital will
care about the cycle-to-cycle jitter affects on setup and hold times.
Radio's will care about the frequency content of the noise with an
offset to the carrier wave.

\subsection{Stability}\label{stability}\index{stability@Stability}

The variation over all corners and enviromental conditions is usually
given in a percentage, parts per million, or parts per billion.

For a digital clock to run a Micro-Controller, maybe it's sufficient
with 10\% accuracy of the clock frequency. For a Bluetooth radio we must
have +-50 ppm, set by the standard. For GPS we might need
parts-per-billion.

\subsection{Conclusion}\label{conclusion}

Each ``clock problem'' will have different frequency, noise and
stability requirements. You must know the order of magnitude of those
before you can design a clock source. There is no ``one-solution fits
all'' clock generation IP.

\section{A typical System-On-Chip clock
system}\label{a-typical-system-on-chip-clock-system}

On the
\href{https://www.nordicsemi.com/Products/Development-hardware/nrf52-dk}{nRF52832
development kit} you can see some components that indicate what type of
clock system must be inside the IC.

In the figure below you can see the following items.

\begin{enumerate}
\def\labelenumi{\arabic{enumi}.}
\tightlist
\item
  32 MHz crystal
\item
  32 KiHz crystal
\item
  PCB antenna
\item
  DC/DC inductor
\end{enumerate}


{
\centering
\aicfig{media/l08_nrf53.png}
}

\emph{Figure 1: nRF5 development kit PCB with (1) 32 MHz crystal, (2) 32
KiHz crystal, (3) PCB antenna, and (4) DC/DC inductor. Source: Nordic
Semiconductor, nRF5340 documentation}

\subsection{32 MHz crystal}\label{mhz-crystal}

Any Bluetooth radio will need a frequency reference. We need to generate
an accurate 2.402 GHz - 2.480 GHz carrier frequency for the gaussian
frequency shift keying (GFSK) modulation. The Bluetooth Standard
requires a +- 50 ppm accurate timing reference, and carrier frequency
offset accuracy.

I'm not sure it's possible yet to make an IC that does not have some
form of frequency reference, like a crystal. The ICs I've seen so far
that have ``crystal less radio'' usually have a resonator (crystal or
bulk-acoustic-wave or MEMS resonator) on die.

The power consumption of a high frequency crystal will be proportional
to frequency. Assuming we have a digital output, then the power of that
digital output will be \(P = C V^2 f\), for example
\(P = 100\text{ fF} \times 1\text{ V}^2 \times 32\text{ MHz} = 3.2\text{ } \mu\text{W}\)
is probably close to a minimum power consumption of a 32 MHz clock.

\subsection{32 KiHz crystal}\label{kihz-crystal}

Reducing the frequency, we can get down to minimum power consumption of
\(P = 100\text{ fF} \times 1\text{ V}^2 \times 32\text{ KiHz} = 3.2 \text{ nW}\)
for a clock.

For a system that sleeps most of the time, and only wakes up at regular
ticks to do something, then a low-frequency crystal might be worth the
effort.

\subsection{PCB antenna}\label{pcb-antenna}\index{pcb antenna@PCB antenna}

Since we can see the PCB antenna, we know that the IC includes a radio.
From that fact we can deduce what must be inside the SoC. If we read the
\href{https://infocenter.nordicsemi.com/index.jsp?topic=\%2Fstruct_nrf52\%2Fstruct\%2Fnrf52832_ps.html}{Product
Specification} we can understand more.

\subsection{DC/DC inductor}\label{dcdc-inductor}\index{dc/dc inductor@DC/DC inductor}

Since we can see a large inductor, we can also make the assumption that
the IC contains a switched regulator. That switched regulator,
especially if it has a pulse-width-modulated control loop, will need a
clock.

From our assumptions we could make a guess what must be inside the IC,
something like the picture below.

There will be a crystal oscillator connected to the crystal. We'll learn
about those later.

These crystal oscillators generate a fixed frequency, 32 MHz, or 32
KiHz, but there might be other clocks needed inside the IC.

To generate those clocks, there will be phase-locked loops (PLL),
frequency locked loops (FLL), or delay-locked loops (DLL).

PLLs take a reference input, and can generate a higher frequency, (or
indeed lower frequency) output. A PLL is a magical block. It's one of
the few analog IPs where we can actually design for infinite gain in our
feedback loop.


{
\centering
\aicfig{media/l10_clockic_tikz.pdf}
}

\emph{Figure 2: Guess at the clock system inside the SoC: crystal
oscillators (XO) for 32 MHz and 32768 Hz, a PLL for the radio local
oscillator, an RC oscillator, and the MCU clock}

Most of the digital blocks on an IC will be synchronous logic, see
figure below. A fundamental principle of synchronous logic is that the
data at the flip-flops (DFF, rectangles with triangle clock input, D, Q
and \(\overline{\text{Q}}\)) only need to be correct at certain times.

The sequence of transitions in the combinatorial logic is of no
consequence, as long as the B inputs are correct when the clock goes
high next time.

The registers, or flip-flops, are your SystemVerilog ``always\_ff''
code. While the blue cloud is your ``always\_comb'' code.

In a SoC we have to check, for all paths between a Y{[}N{]} and B{[}M{]}
that the path is fast enough for all transients to settle before the
clock strikes next time. How early the B data must arrive in relation to
the clock edge is the setup time of the DFFs.

We also must check for all paths that the B{[}M{]} are held for long
enough after the clock strikes such that our flip-flop does not change
state. The hold time is the distance from the clock edge to where the
data is allowed to change. Negative hold times are common in DFFs, so
the data can start to change before the clock edge.

In an IC with millions of flip-flops there can be billions of paths. The
setup and hold time for every single one must be checked. One could
imagine a simulation of all the paths on a netlist with parasitics
(capacitors and resistors from layout) to check the delays, but there
are so many combinations that the simulation time becomes unpractical.

Static Timing Analysis (STA) is a light-weight way to check all the
paths. For the STA we make a model of the delay in each cell (captured
in a liberty file), the setup/hold times of all flip-flops, wire
propagation delays, clock frequency (or period), and the variation in
the clock frequency. The process, voltage, temperature variation must
also be checked for all components, so the number of liberty files can
quickly grow large.

For an analog designer the constraints from digital will tell us what's
the maximum frequency we can have at any point in time, and what is the
maximum cycle-to-cycle variation in the period.


{
\centering
\aicfig{media/logic_tikz.pdf}
}

\emph{Figure 3: Synchronous logic: flip-flops capture the data on the
clock edge, with combinatorial logic (blue cloud) between register
stages}

\section{PLL}\label{pll}\index{pll@PLL}

PLL, or it's cousins FLL and DLL are really cool. A PLL is based on the
familiar concept of feedback, shown in the figure below. As long as we
make \(H(s)\) infinite we can force the output to be an exact copy of
the input.


{
\centering
\aicfig{media/l10_fb_tikz.pdf}
}

\emph{Figure 4: Feedback loop where an infinite gain H(s) forces the
output to be an exact copy of the input}

\subsection{Integer PLL}\label{integer-pll}\index{integer pll@Integer PLL}

For a frequency loop the figure looks a bit different. If we want a
higher output frequency we can divide the frequency by a number (N) and
compare with our reference (for example the 32 MHz reference from the
crystal oscillator).

We then take the error, apply a transfer function \(H(s)\) with high
gain, and control our oscillator frequency.

If the down-divided output frequency is too high, we force the
oscillator to a lower frequency. If the down-divided output frequency is
too low we force the oscillator to a higher frequency.

If we design the \(H(s)\) correctly, then we have
\(f_o = N \times f_{in}\)


{
\centering
\aicfig{media/l10_freq_fb_tikz.pdf}
}

\emph{Figure 5: Integer PLL: the oscillator output is divided by N and
compared to the reference, giving an output frequency N times the
reference}

Sometimes you want a finer frequency resolution, in that case you'd add
a divider on the reference and get \(f_o = N \times \frac{f_{in}}{M}\)..


{
\centering
\aicfig{media/l08_pll_m_tikz.pdf}
}

\emph{Figure 6: Integer PLL with an additional divide-by-M on the
reference for finer frequency resolution}

\subsection{Fractional PLL}\label{fractional-pll}\index{fractional pll@Fractional PLL}

Trouble is that dividing down the input frequency will reduce your loop
bandwidth, as the low-pass filter needs to be about 1/10'th of the
reference frequency. As such, the PLL will respond slower to a frequency
change.

We can also use a fractional divider, where we swap between two, or
more, integers in a sigma-delta fashion in the divider.


{
\centering
\aicfig{media/l08_pll_sd_tikz.pdf}
}

\emph{Figure 7: Fractional PLL where a sigma-delta modulator switches
the feedback divider between integer values}

\subsection{Modulation in PLLs}\label{modulation-in-plls}\index{modulation in plls@Modulation in PLLs}

From your signal processing, or communication courses, you may recognize
the equation below.

\[ A_m(t) \times cos\left( 2 \pi f_{carrier}t + \phi_{m}(t)\right)\]

The \(A_m\) is the amplitude modulation, while the \(\phi_m\) is the
phase modulation. Bluetooth Low Energy is constant envelope, so the
\(A_m\) is a constant. The phase modulation is applied to the carrier,
but how is it done?

One option is shown below. We could modulate our frequency reference
directly. That could maybe be a sigma-delta divider on the reference, or
directly modulating the oscillator.


{
\centering
\aicfig{media/l08_pll_mod_tikz.pdf}
}

\emph{Figure 8: Modulating the PLL by adding the modulation signal
directly to the frequency reference}

Most modern radios, however, will have a two-point modulation. The
modulation signal is applied to the VCO (or DCO), and the opposite
signal is applied to the feedback divider. As such, the modulation is
not seen by the loop.


{
\centering
\aicfig{media/l08_pll_2mod_tikz.pdf}
}

\emph{Figure 9: Two-point modulation: the modulation is applied to the
oscillator and the opposite signal to the sigma-delta feedback divider,
so the loop does not see it}

\section{PLL Example}\label{pll-example}\index{pll example@PLL Example}

I've made an example
\href{https://github.com/wulffern/sun_pll_sky130nm}{PLL} that you can
download and play with. I make no claims that it's a good PLL. Actually,
I know it's a bad PLL. The ring-oscillator frequency varies too fast
with the voltage control. But it does give you a starting point.

A PLL can consist of a oscillator (SUN\_PLL\_ROSC) that generates our
output frequency. A divider (SUN\_PLL\_DIVN) that generates a feedback
frequency that we can compare to the reference. A Phase and Frequency
Detector (SUN\_PLL\_PFD) and a charge-pump (SUN\_PLL\_CP) that model the
\(+\), or the comparison function in our previous picture. And a loop
filter (SUN\_PLL\_LPF and SUN\_PLL\_BUF) that is our \(H(s)\).


{
\centering
\aicfig{media/sunpll_top_tikz.pdf}
}

\emph{Figure 10: Top-level schematic of the SUN\_PLL example with
phase-frequency detector, charge-pump, loop filter, buffer, ring
oscillator, and divide-by-32 feedback divider. Block placement follows
the xschem source: signal flow left to right along the loop, feedback
below, bias and start-up at the bottom}

Read any book on PLLs, talk to any PLL designer and they will all tell
you the same thing. \textbf{PLLs require calculation}. You must setup a
linear model of the feedback loop, and calculate the loop transfer
function to check the stability, and the loop gain. \textbf{This is the
way!} (to quote Mandalorian).

But how can we make a linear model of a non-linear system? The voltages
inside a PLL must be non-linear, they are clocks. A PLL is not linear in
time-domain!

I have no idea who first thought of the idea, but it turns out, that one
can model a PLL as a linear system if one consider the phase of the
voltages inside the PLL, especially when the PLL is locked (phase of the
output and reference is mostly aligned). Where the phase is defined as

\[ \phi(t) = 2 \pi \int_0^t f(\tau) d\tau\]

As long as the bandwidth of the \(H(s)\) is about \(\frac{1}{10}\) of
the reference frequency, then the linear model below holds (at least is
good enough).

The phase of our input is \(\phi_{in}(s)\), the phase of the output is
\(\phi(s)\), the divided phase is \(\phi_{div}(s)\) and the phase error
is \(\phi_d(s)\).

The \(K_{pd}\) is the gain of our phase-frequency detector and
charge-pump. The \(K_{lp}H_{lp}(s)\) is our loop filter \(H(s)\). The
\(K_{osc}/s\) is our oscillator transfer function. And the \(1/N\) is
our feedback divider.


{
\centering
\aicfig{media/l10_pll_sm_tikz.pdf}
}

\emph{Figure 11: Linear phase-domain model of the PLL with
phase-detector gain, loop filter, oscillator integrator and 1/N feedback
divider}

\subsection{Loop gain}\label{loop-gain}\index{loop gain@Loop gain}

The loop transfer function can then be analyzed and we get.

\[ \frac{\phi_d}{\phi_{in}} = \frac{1}{1 + L(s)}\]

\[ L(s) = \frac{ K_{osc} K_{pd} K_{lp} H_{lp}(s) }{N s} \]

Here is the magic of PLLs. Notice what happens when
\(s = j\omega = j 0\), or at zero frequency. If we assume that
\(H_{lp}(s)\) is a low pass filter, then
\(H_{lp}(0) = \text{constant}\). The loop gain, however, will have a
\(L(0) \propto \frac{1}{0}\) which approaches infinity at 0.

That means, we have an infinite DC gain in the loop transfer function.
It is the only case I know of in an analog design where we can actually
have infinite gain. Infinite gain translates to infinite precision.

If the reference was a Rubidium oscillator we could generate any
frequency with the same precision as the frequency of the Rubidium
oscillator. Magic.

For the linear model, we need to figure out the factors, like
\(K_{osc}\), which must be determined by simulation.

\subsection{Controlled oscillator}\label{controlled-oscillator}\index{controlled oscillator@Controlled oscillator}

The gain of the oscillator is the change in output frequency as a
function of the change of the control node. For a voltage-controlled
oscillator (VCO) we could sweep the control voltage, and check the
frequency. The derivative of the f(V) would be proportional to the
\(K_{vco}\).

The control node does not need to be a voltage. Anything that changes
the frequency of the oscillator can be used as a control node. There
exist PLLs with voltage control, current control, capacitance control,
and digital control.

For the SUN\_PLL\_ROSC it is the VDD of the ring-oscillator (VDD\_ROSC)
that is our control node.

\[K_{osc} = 2 \pi\frac{ df}{dV_{cntl}}\]


{
\centering
\aicfig{media/sunpll_rosc_tikz.pdf}
}

\emph{Figure 12: The ring oscillator SUN\_PLL\_ROSC: a NAND and eight
inverters make nine inversions, so the loop oscillates whenever
PWRUP\_1V8 is high. The ring runs on VDD\_ROSC --- the supply is the
control node --- and the level shifter LS brings taps N2/N1 back to the
AVDD domain to make CK}

\subsubsection{\texorpdfstring{\href{https://github.com/wulffern/sun_pll_sky130nm/tree/main/sim/ROSC}{SUN\_PLL\_SKY130NM/sim/ROSC/}}{SUN\_PLL\_SKY130NM/sim/ROSC/}}\label{sun_pll_sky130nmsimrosc}

I simulate the ring oscillator in ngspice with a transient simulation
and get the oscillator frequency as a function of voltage.

\textbf{tran.spi}

\begin{Shaded}
\begin{Highlighting}[]
\NormalTok{let start\_v = 1.1}
\NormalTok{let stop\_v = 1.7}
\NormalTok{let delta\_v = 0.1}
\NormalTok{let v\_act = start\_v}
\NormalTok{* loop}
\NormalTok{while v\_act le stop\_v}
\NormalTok{alter VROSC v\_act}
\NormalTok{tran 1p 40n}
\NormalTok{meas tran vrosc avg v(VDD\_ROSC)}
\NormalTok{meas tran tpd trig v(CK) val=\textquotesingle{}0.8\textquotesingle{} rise=10 targ v(CK) val=\textquotesingle{}0.8\textquotesingle{} rise=11}
\NormalTok{let v\_act = v\_act + delta\_v}
\NormalTok{end}
\end{Highlighting}
\end{Shaded}

I use \texttt{tran.py} to extract the time-domain signal from ngspice
into a CSV file.

Then I use a python script to extract the \(K_{osc}\)

\textbf{kvco.py}

\begin{Shaded}
\begin{Highlighting}[]
\NormalTok{    df }\OperatorTok{=}\NormalTok{ pd.read\_csv(f)}
\NormalTok{    freq }\OperatorTok{=} \DecValTok{1}\OperatorTok{/}\NormalTok{df[}\StringTok{"tpd"}\NormalTok{]}
\NormalTok{    kvco }\OperatorTok{=}\NormalTok{ np.mean(freq.diff()}\OperatorTok{/}\NormalTok{df[}\StringTok{"vrosc"}\NormalTok{].diff())}
\end{Highlighting}
\end{Shaded}

Below I've made a plot of the oscillation frequency over corners.


{
\centering
\aicfig{media/SUN_PLL_ROSC_KVCO_tikz.pdf}
}

\emph{Figure 13: Ring-oscillator frequency versus control voltage
VDD\_ROSC over nine process and temperature corners, simulated on the
extracted layout. The slope at the typical corner is 1.01 GHz/V, which
is \(K_{osc}\); the spread is a factor of eleven at 1.2 V, narrowing to
under three at 1.5 V}

Two things in that plot are worth more than the slope.

The first is the spread. A ring oscillator has nothing setting its
frequency except how fast its own inverters switch, so process and
temperature move it by a factor of eleven at the bottom of the control
range. Every corner does cross the 256 MHz the loop needs, but the
slow-cold one only at about 1.44 V, near the top of what the control
node can deliver. The usable tuning range is not the width of the
control range, it is whatever is left above that crossing in the worst
corner.

The second is the missing point. At slow-cold and 1.1 V there is no
measurement, because the oscillator was too slow to produce enough edges
inside the simulated window and the measurement failed rather than
returning a plausible wrong number. A failed measurement in the corner
you were already worried about is information, not an inconvenience ---
and it is a good argument for reading the simulator's errors rather than
only its plots.

\subsection{Phase detector and charge
pump}\label{phase-detector-and-charge-pump}

The two blocks compare our reference clock to our feedback clock, and
produce an error signal. The gain of the pair is the average current fed
into the loop filter per radian of phase error, and it is worth deriving
once because the \(2\pi\) looks arbitrary until you do.

The phase-frequency detector turns a phase error into a pulse width. If
the feedback clock arrives late by a phase \(\Delta\phi\), the UP output
is high for the fraction \(\Delta\phi/2\pi\) of the reference period,
because a full period is \(2\pi\) of phase. During that pulse the charge
pump sources its full current \(I_{cp}\), and for the rest of the period
it sources nothing. The average current into the filter is therefore

\[ \overline{I} = I_{cp}\frac{\Delta\phi}{2\pi} \]

and the gain, being average current per radian, is what is left when you
divide by \(\Delta\phi\):

\[ K_{pd} = \frac{I_{cp}}{2 \pi} \]

Two things follow that are easy to miss. The gain does not depend on the
reference frequency, because both the pulse width and the period scale
together. And it is the \emph{average} current that the loop filter
sees, which is only a fair description if the filter is slow compared
with the reference --- the same assumption that let us draw a linear
model in the first place.


{
\centering
\aicfig{media/sunpll_pfd_tikz.pdf}
}

\emph{Figure 14: The phase-frequency detector SUN\_PLL\_PFD: two
flip-flops with D tied high, one set by CK\_REF, the other by CK\_FB.
The moment both are set the NOR resets the pair, so the surviving pulse
width is the arrival-time difference --- phase error becomes pulse
width}


{
\centering
\aicfig{media/sunpll_cp_tikz.pdf}
}

\emph{Figure 15: The charge pump SUN\_PLL\_CP, driven by the
phase-frequency detector through CP\_UP\_N and CP\_DOWN. V\_BN sets
I\_cp, the switch pair steers it into or out of V\_LPF, M7 parks V\_LPF
at AVDD in power-down, and the KICK switch grabs the filter's zero node
to start the loop. The mirror devices are stacked pairs in silicon,
drawn single here}

\subsection{Loop filter}\label{loop-filter}\index{loop filter@Loop filter}

In the book you'll find a first order loop filter, and a second order
loop filter. Engineers are creative, so you'll likely find other loop
filters in the literature.

I would start with the ``known to work'' loop filters before you explore
on your own.

If you're really interested in PLLs, you should buy
\href{https://www.amazon.com/Design-CMOS-Phase-Locked-Loops-Architecture/dp/1108494544}{Design
of CMOS Phase-Locked Loops} by Behzad Razavi.

The loop filter has a unity gain buffer. My oscillator draws current,
while the VLPF node is high impedance, so I can't draw current from the
loop filter without changing the filter transfer function.

\[ K_{lp}H_{lp}(s)= K_{lp}\left(\frac{1}{s} + \frac{1}{\omega_z}\right) \]

\[ K_{lp}H_{lp}(s) = \frac{1}{s(C_1 + C_2)}\frac{1 + s R C_1}{1 +
sR\frac{C_1C_2}{C_1 + C_2}}\]


{
\centering
\aicfig{media/sunpll_lpf_tikz.pdf}
}

\emph{Figure 16: The loop filter SUN\_PLL\_LPF on VLPF, followed by the
buffer SUN\_PLL\_BUF that drives the oscillator supply VDD\_ROSC. C1 is
22 unit capacitors, C2 is 3, so the zero sits where the equations above
put it}

\subsection{Divider}\label{divider}\index{divider@Divider}

The divider is modelled as

\[ K_{div} = \frac{1}{N}\]


{
\centering
\aicfig{media/sunpll_divn_tikz.pdf}
}

\emph{Figure 17: The feedback divider SUN\_PLL\_DIVN: five flip-flops
wired as toggles, each clocking the next, dividing CK by 32 to make
CK\_FB}

\subsection{Loop transfer function}\label{loop-transfer-function}\index{loop transfer function@Loop transfer function}

With the loop transfer function we can start to model what happens in
the linear loop. What is the phase response, and what is the gain
response.

\[ L(s) = \frac{ K_{osc} K_{pd} K_{lp} H_{lp}(s) }{N s} \]

\subsubsection{Python model}\label{python-model}

I've made a python model of the loop, you can find it at
\href{https://github.com/wulffern/sun_pll_sky130nm/blob/main/jupyter/pll.ipynb}{sun\_pll\_sky130nm/jupyter/pll}
-
\href{https://wulffern.github.io/aic2026/assets/examples/pll.html}{interactive}

In the jupyter notebook below you can find some more information on the
phase/frequency detector, and charge pump.

\href{https://github.com/wulffern/sun_pll_sky130nm/blob/main/jupyter/pfd.ipynb}{sun\_pll\_sky130nm/jupyter/pfd}

Below is a plot of the loop gain, and the transfer function from input
phase to divider phase.

We can see that the loop gain at low frequency is large, and
proportional to \(1/s\). As such, the phase of the divided down feedback
clock is the same as our reference.

The closed loop transfer function \(\phi_{div}/\phi_{in}\) shows us that
the divided phase at low frequency is the same as the input phase. Since
the phase is the same, and the frequency must be the same, then we know
that the output clock will be N times reference frequency.

Which \(K_{osc}\) went into that plot matters more than it looks. The
schematic simulation gave 1.6 GHz/V; the extracted layout gives 1.01
GHz/V, and the difference is parasitic capacitance in the ring that
simply does not exist until the oscillator is laid out. A third of the
gain disappears, and since the loop gain is proportional to \(K_{osc}\),
the crossover moves from 0.59 MHz down to 0.43 MHz and the phase margin
from 51 degrees to 43.

Eight degrees is not a catastrophe, and that is rather the point: it is
the sort of erosion that is easy to spend twice over without noticing.
Extract early, and design the loop with the number the silicon will
actually have rather than the one the schematic promised.

It is also worth checking this plot against the assumption we made when
we drew the linear model at all. The loop gain crosses 0 dB at 0.43 MHz,
and the reference frequency is \(256\ \text{MHz}/32 = 8\) MHz, so the
loop bandwidth is a nineteenth of the reference. That clears the ``one
tenth of the reference'' rule, which means the model is entitled to be
believed. If it had not cleared it, the phase margin the plot reports
would be a number about a model that does not describe the circuit ---
and that is a far worse situation than a poor phase margin, because it
looks fine.


{
\centering
\aicfig{media/pll_tikz.pdf}
}

\emph{Figure 18: Magnitude and phase of the loop gain and the
closed-loop transfer function from input phase to divider phase, using
the oscillator gain measured on the extracted layout. The loop crosses 0
dB at 0.43 MHz with 43 degrees of phase margin}

The top testbench for the PLL is
\href{https://github.com/wulffern/sun_pll_sky130nm/blob/main/sim/SUN_PLL/tran.spi}{tran.spi}.

I power up the PLL and wait for the output clock to settle. The
frequency is measured the way a counter would measure it: find every
rising edge of CK and take the reciprocal of the interval between
consecutive edges. See
\href{https://github.com/wulffern/sun_pll_sky130nm/blob/main/sim/SUN_PLL/freq.py}{freq.py}.


{
\centering
\aicfig{media/sun_pll_lay_typ_tikz.pdf}
}

\emph{Figure 19: Simulated PLL output frequency from power-up, on the
extracted layout at the typical corner. The grey trace is the frequency
of each individual cycle and the black one a 200 cycle average; the loop
overshoots to about 500 MHz, undershoots past the target, and settles at
256.1 MHz around 12 microseconds}

Three things in that plot are worth pausing on.

The loop starts at the top of the oscillator's range and has to be
dragged down, so the first microsecond is not feedback at all, it is the
control node charging. Then the loop takes over, overshoots past the
target, and rings once before settling. That single undershoot is the
second order response the phase margin describes; 43 degrees is what one
visible ring looks like.

The grey band does not narrow as the loop settles. That is not the
simulation failing to converge, and it is not an error the loop could
correct: the charge pump delivers its correction as a pulse once per
reference period, so the oscillator is kicked 8 million times a second
whatever the loop is doing. In a real chip this is the reference spur,
and it is the reason a PLL's output is never as clean as its reference.

Settling takes about 12 microseconds here, against roughly 8 in earlier
versions of this design. The loop got slower because the oscillator got
slower: the extracted layout has a third less gain than the schematic,
and loop bandwidth is proportional to that gain. Nothing was designed
differently; the parasitics simply arrived.

You can find the schematics, layout, testbenches, python script etc at
\href{https://github.com/wulffern/sun_pll_sky130nm}{SUN\_PLL\_SKY130NM}

Below are a couple layout images of the finished PLL


{
\centering
\aicfig{media/sun_pll_layout0.png}
}

\emph{Figure 20: Floorplan of the SUN\_PLL layout; the loop filter
capacitor SUN\_PLL\_LPF dominates the area above the PLL blocks}


{
\centering
\aicfig{media/sun_pll_layout1.png}
}

\emph{Figure 21: Finished SUN\_PLL layout showing the loop filter
capacitor array and the PLL blocks along the bottom}

\section{Summary}\label{summary}

The one-page version of this chapter:

\begin{itemize}
\tightlist
\item
  Everything on the chip wants a clock, and every clock is a compromise
  between frequency accuracy (ppm), phase noise and power
\item
  A crystal gives the accurate reference; the PLL multiplies it up to
  the frequency the system needs
\item
  PFD turns phase error into pulse width, the charge pump into charge,
  the loop filter into a control voltage, the oscillator into frequency,
  and the divider closes the loop
\item
  The type-II loop needs its zero: C1 sets the zero with R, C2 cleans
  the ripple, and the loop bandwidth balances reference noise against
  VCO noise
\item
  Inside the bandwidth the PLL follows the reference; outside, the
  oscillator is on its own - phase noise plots read exactly that way
\item
  SUN\_PLL is the whole story in five schematics: ROSC, PFD, charge
  pump, loop filter, divider
\end{itemize}

\section{Would you like to know
more?}\label{would-you-like-to-know-more}

Back in 2020 there was a Master student at NTNU on PLL. I would
recommend looking at that thesis to learn more, and to get inspired
\href{https://ntnuopen.ntnu.no/ntnu-xmlui/handle/11250/2778127}{Ultra
Low Power Frequency Synthesizer}.

A Low Noise Sub-Sampling PLL in Which Divider Noise is Eliminated and
PD/CP Noise is Not Multiplied by N2 \citeproc{ref-gao09}{{[}1{]}}

All-digital PLL and transmitter for mobile phones
\citeproc{ref-staszewski05}{{[}2{]}}

A 2.9--4.0-GHz Fractional-N Digital PLL With Bang-Bang Phase Detector
and 560-fsrms Integrated Jitter at 4.5-mW Power
\citeproc{ref-tasca11}{{[}3{]}}

\phantomsection\label{refs}
\begin{CSLReferences}{0}{0}
\bibitem[\citeproctext]{ref-gao09}
\CSLLeftMargin{{[}1{]} }%
\CSLRightInline{X. Gao, E. A. M. Klumperink, M. Bohsali, and B. Nauta,
{``A low noise sub-sampling PLL in which divider noise is eliminated and
PD/CP noise is not multiplied by \(N ^{2}\),''} \emph{IEEE Journal of
Solid-State Circuits}, vol. 44, no. 12, pp. 3253--3263, 2009, doi:
\href{https://doi.org/10.1109/JSSC.2009.2032723}{10.1109/JSSC.2009.2032723}.
}

\bibitem[\citeproctext]{ref-staszewski05}
\CSLLeftMargin{{[}2{]} }%
\CSLRightInline{R. Staszewski \emph{et al.}, {``All-digital PLL and
transmitter for mobile phones,''} \emph{IEEE Journal of Solid-State
Circuits}, vol. 40, no. 12, pp. 2469--2482, 2005, doi:
\href{https://doi.org/10.1109/JSSC.2005.857417}{10.1109/JSSC.2005.857417}.
}

\bibitem[\citeproctext]{ref-tasca11}
\CSLLeftMargin{{[}3{]} }%
\CSLRightInline{D. Tasca, M. Zanuso, G. Marzin, S. Levantino, C. Samori,
and A. L. Lacaita, {``A 2.9--4.0-GHz fractional-n digital PLL with
bang-bang phase detector and 560-\(\mathrm{fs}_\mathrm{rms}\) integrated
jitter at 4.5-mW power,''} \emph{IEEE Journal of Solid-State Circuits},
vol. 46, no. 12, pp. 2745--2758, 2011, doi:
\href{https://doi.org/10.1109/JSSC.2011.2162917}{10.1109/JSSC.2011.2162917}.
}

\end{CSLReferences}
